% Options for packages loaded elsewhere
% Options for packages loaded elsewhere
\PassOptionsToPackage{unicode}{hyperref}
\PassOptionsToPackage{hyphens}{url}
\PassOptionsToPackage{dvipsnames,svgnames,x11names}{xcolor}
%
\documentclass[
  11pt,
]{article}
\usepackage{xcolor}
\usepackage[margin=1in]{geometry}
\usepackage{amsmath,amssymb}
\setcounter{secnumdepth}{-\maxdimen} % remove section numbering
\usepackage{iftex}
\ifPDFTeX
  \usepackage[T1]{fontenc}
  \usepackage[utf8]{inputenc}
  \usepackage{textcomp} % provide euro and other symbols
\else % if luatex or xetex
  \usepackage{unicode-math} % this also loads fontspec
  \defaultfontfeatures{Scale=MatchLowercase}
  \defaultfontfeatures[\rmfamily]{Ligatures=TeX,Scale=1}
\fi
\usepackage{lmodern}
\ifPDFTeX\else
  % xetex/luatex font selection
\fi
% Use upquote if available, for straight quotes in verbatim environments
\IfFileExists{upquote.sty}{\usepackage{upquote}}{}
\IfFileExists{microtype.sty}{% use microtype if available
  \usepackage[]{microtype}
  \UseMicrotypeSet[protrusion]{basicmath} % disable protrusion for tt fonts
}{}
\makeatletter
\@ifundefined{KOMAClassName}{% if non-KOMA class
  \IfFileExists{parskip.sty}{%
    \usepackage{parskip}
  }{% else
    \setlength{\parindent}{0pt}
    \setlength{\parskip}{6pt plus 2pt minus 1pt}}
}{% if KOMA class
  \KOMAoptions{parskip=half}}
\makeatother
% Make \paragraph and \subparagraph free-standing
\makeatletter
\ifx\paragraph\undefined\else
  \let\oldparagraph\paragraph
  \renewcommand{\paragraph}{
    \@ifstar
      \xxxParagraphStar
      \xxxParagraphNoStar
  }
  \newcommand{\xxxParagraphStar}[1]{\oldparagraph*{#1}\mbox{}}
  \newcommand{\xxxParagraphNoStar}[1]{\oldparagraph{#1}\mbox{}}
\fi
\ifx\subparagraph\undefined\else
  \let\oldsubparagraph\subparagraph
  \renewcommand{\subparagraph}{
    \@ifstar
      \xxxSubParagraphStar
      \xxxSubParagraphNoStar
  }
  \newcommand{\xxxSubParagraphStar}[1]{\oldsubparagraph*{#1}\mbox{}}
  \newcommand{\xxxSubParagraphNoStar}[1]{\oldsubparagraph{#1}\mbox{}}
\fi
\makeatother


\usepackage{longtable,booktabs,array}
\newcounter{none} % for unnumbered tables
\usepackage{calc} % for calculating minipage widths
% Correct order of tables after \paragraph or \subparagraph
\usepackage{etoolbox}
\makeatletter
\patchcmd\longtable{\par}{\if@noskipsec\mbox{}\fi\par}{}{}
\makeatother
% Allow footnotes in longtable head/foot
\IfFileExists{footnotehyper.sty}{\usepackage{footnotehyper}}{\usepackage{footnote}}
\makesavenoteenv{longtable}
\usepackage{graphicx}
\makeatletter
\newsavebox\pandoc@box
\newcommand*\pandocbounded[1]{% scales image to fit in text height/width
  \sbox\pandoc@box{#1}%
  \Gscale@div\@tempa{\textheight}{\dimexpr\ht\pandoc@box+\dp\pandoc@box\relax}%
  \Gscale@div\@tempb{\linewidth}{\wd\pandoc@box}%
  \ifdim\@tempb\p@<\@tempa\p@\let\@tempa\@tempb\fi% select the smaller of both
  \ifdim\@tempa\p@<\p@\scalebox{\@tempa}{\usebox\pandoc@box}%
  \else\usebox{\pandoc@box}%
  \fi%
}
% Set default figure placement to htbp
\def\fps@figure{htbp}
\makeatother





\setlength{\emergencystretch}{3em} % prevent overfull lines

\providecommand{\tightlist}{%
  \setlength{\itemsep}{0pt}\setlength{\parskip}{0pt}}



 


\makeatletter
\@ifpackageloaded{caption}{}{\usepackage{caption}}
\AtBeginDocument{%
\ifdefined\contentsname
  \renewcommand*\contentsname{Table of contents}
\else
  \newcommand\contentsname{Table of contents}
\fi
\ifdefined\listfigurename
  \renewcommand*\listfigurename{List of Figures}
\else
  \newcommand\listfigurename{List of Figures}
\fi
\ifdefined\listtablename
  \renewcommand*\listtablename{List of Tables}
\else
  \newcommand\listtablename{List of Tables}
\fi
\ifdefined\figurename
  \renewcommand*\figurename{Figure}
\else
  \newcommand\figurename{Figure}
\fi
\ifdefined\tablename
  \renewcommand*\tablename{Table}
\else
  \newcommand\tablename{Table}
\fi
}
\@ifpackageloaded{float}{}{\usepackage{float}}
\floatstyle{ruled}
\@ifundefined{c@chapter}{\newfloat{codelisting}{h}{lop}}{\newfloat{codelisting}{h}{lop}[chapter]}
\floatname{codelisting}{Listing}
\newcommand*\listoflistings{\listof{codelisting}{List of Listings}}
\makeatother
\makeatletter
\makeatother
\makeatletter
\@ifpackageloaded{caption}{}{\usepackage{caption}}
\@ifpackageloaded{subcaption}{}{\usepackage{subcaption}}
\makeatother
\usepackage{bookmark}
\IfFileExists{xurl.sty}{\usepackage{xurl}}{} % add URL line breaks if available
\urlstyle{same}
\hypersetup{
  pdftitle={Ice Chamber Experiment Metagenome},
  pdfkeywords={microbial ecology, metagenomics},
  colorlinks=true,
  linkcolor={blue},
  filecolor={Maroon},
  citecolor={Blue},
  urlcolor={Blue},
  pdfcreator={LaTeX via pandoc}}


\title{Ice Chamber Experiment Metagenome}
\author{}
\date{2026-07-22}
\begin{document}
\maketitle
\begin{abstract}
Sea ice and its associated frost flowers host taxonomically and
functionally distinct microbial communities that play important roles in
polar biogeochemical cycles, yet field-based studies of these systems
are complicated by confounding regional geography and large-scale ocean
and atmospheric circulation, which obscure the local drivers of
microbial colonization (Priest et al., 2023). To address this
limitation, we used the Roland von Glasow Air-Sea-Ice Chamber, a large
laboratory-scale sea ice mesocosm at the University of Grenoble Alpes,
to characterize bacterial and eukaryotic community assembly across frost
flower, brine, ice matrix, and seawater niches under controlled
conditions. A total of 124 SRA runs from 45 samples were processed using
an amplicon sequence variant (ASV)-based bioinformatics pipeline,
including quality filtering, taxonomic classification, and functional
annotation. Of 84 samples subjected to ASV inference, an initial 860,222
reads were quality filtered to 519,769 retained reads (60.4\%
retention). Analysis of the resulting community profiles revealed
taxonomic and functional partitioning of bacteria into distinct niches
consistent with niche-based selection, supported by the emergence of
archetypal sea-ice bacterial taxa that were not detected in the seeding
seawater. In contrast, eukaryotic community composition remained largely
consistent across the ice-brine-seawater profile, suggesting that
eukaryotic distribution in this system was governed more by stochastic
processes than by deterministic niche selection. Bacterial enrichment
within the ice matrix proceeded at a rate comparable to that reported in
prior field-based observations, despite minimal diatom representation in
this experimental system. These findings demonstrate that controlled
mesocosm experiments can effectively isolate local physicochemical
drivers of microbial colonization from confounding environmental
variability, and support a model in which bacterial and eukaryotic
communities in sea ice are shaped by distinct assembly
processes---niche-based selection and ecological stochasticity,
respectively.
\end{abstract}


\subsection{Introduction}\label{introduction}

\section{Introduction}\label{introduction-1}

Sea ice and its associated frost flowers---delicate, salt-rich
structures that form on the surface of newly frozen sea ice---host
diverse microbial communities that contribute meaningfully to polar
biogeochemical cycles, including carbon flux, nutrient recycling, and
the production of climate-active aerosols {[}ref15{]}. Despite their
ecological importance, these communities are challenging to study in
natural settings, where large-scale atmospheric and oceanic dynamics
confound efforts to isolate the local, deterministic forces governing
microbial colonization and assembly {[}ref14{]}. Distinguishing
niche-based selection from stochastic dispersal processes is a central
question in microbial ecology {[}ref13{]}, yet field studies of sea ice
are inherently limited in their ability to control for the many
interacting physical and chemical gradients that develop during ice
formation. Laboratory-scale simulations offer an alternative: the Roland
von Glasow Air-Sea-Ice Chamber, a large-scale experimental facility
capable of generating young sea ice and frost flowers under controlled
conditions, provides an opportunity to examine microbial colonization
while minimizing the confounding influence of regional atmospheric and
oceanic variability.

In this study, we used the Air-Sea-Ice Chamber at the University of
Grenoble-Alpes to investigate how bacterial and eukaryotic assemblages
are structured across the distinct physicochemical niches that develop
within experimentally generated sea ice---frost flowers, brine, the ice
matrix, and the underlying seawater. Building on prior observations that
bacterial taxa in sea ice environments often display strong
niche-specific enrichment {[}ref12{]}, we hypothesized that bacterial
communities in the chamber system would exhibit significant taxonomic
partitioning among these habitats, consistent with deterministic
niche-based selection. We further predicted that this partitioning would
be evidenced by the emergence of archetypal sea-ice bacterial taxa that
were not present in the seawater used to seed the experiment. By
contrast, we anticipated that eukaryotic assemblages would show
comparatively little variation across the ice profile, reflecting a
greater role for stochastic colonization processes among these taxa.
Finally, we expected that enrichment of characteristic sea-ice bacteria
would be detectable even under conditions of minimal diatom
representation, which would underscore the robustness of bacterial
responses to ice-associated physicochemical conditions independent of
eukaryotic primary production.

To test these hypotheses, we reanalyzed publicly available 16S rRNA gene
amplicon sequencing data generated from the Air-Sea-Ice Chamber
experiment (BioProject PRJNA1473294), comprising 45 biological samples
and 124 sequencing runs spanning the frost flower, brine, ice matrix,
and seawater niches. Raw sequence reads were processed using the
microscape pipeline, which incorporates quality filtering, primer
removal, denoising via DADA2, chimera removal, taxonomic classification
against the SILVA reference database, and decontamination of likely
contaminant sequences. The resulting amplicon sequence variant table was
used to characterize taxonomic composition across niches and to evaluate
patterns of alpha- and beta-diversity, differential taxon abundance, and
co-occurrence structure. This approach allowed us to assess, at high
resolution, the extent to which bacterial and eukaryotic community
structure in this experimentally controlled sea-ice system reflects
deterministic niche selection versus stochastic assembly processes.

\subsection{Methods}\label{methods}

\section{Methods}\label{methods-1}

\subsection{Data Acquisition}\label{data-acquisition}

Sequence data analyzed in this study were derived from an experimental
sea ice mesocosm study conducted using the Roland von Glasow Air-Sea-Ice
Chamber (University of Grenoble Alpes), designed to isolate local
microbial colonization dynamics on frost flowers, brine, ice matrix, and
seawater from confounding regional environmental effects. A total of 45
samples were processed, corresponding to 124 sequencing runs deposited
under BioProject PRJNA1473294 in the NCBI Sequence Read Archive (SRA).
Samples represented air metagenome and three additional sample
types/environments consistent with the frost flower, brine, ice matrix,
and seawater niches described in the study design {[}AUTHOR: confirm
sample-to-environment mapping and per-environment sample counts{]}. Raw
sequencing reads were retrieved from SRA using standard SRA toolkit
utilities prior to downstream processing.

\subsection{Sequence Processing
Pipeline}\label{sequence-processing-pipeline}

Amplicon sequence data were processed using the microscape pipeline,
which implements an amplicon sequence variant (ASV)-based workflow
analogous to DADA2-style denoising. Following quality filtering, an
initial set of 84 samples with sufficient read depth were retained for
downstream community analysis (out of the total 45 biological
samples/124 sequencing runs, indicating that multiple runs were merged
or samples were processed at the run level for this step {[}AUTHOR:
clarify sample-to-run aggregation strategy{]}). Quality filtering
removed low-quality and chimeric reads, reducing the initial 860,222
input reads to 519,769 retained reads, corresponding to an overall
retention rate of 60.4\%.

Filtered reads were used to construct sequence tables (seqtabs) and,
following merging and chimera removal, a final ASV table (seqtab\_final)
was generated. Taxonomic classification of ASVs was performed against
the SILVA reference database, which is the sole taxonomy database
employed in this analysis (version not specified in pipeline metadata)
{[}AUTHOR: specify SILVA release version if known{]}. Classification was
conducted hierarchically across taxonomic ranks from domain to genus.

\subsection{Analysis Parameters}\label{analysis-parameters}

A total of 162 ASVs were classified using the SILVA database.
Classification success was high and consistent across upper taxonomic
ranks, with all 162 ASVs (100\%) assigned at the domain, phylum, and
class levels. Assignment rates declined marginally at finer resolution,
with 161 ASVs (99.4\%) classified at order and family levels, and 152
ASVs (93.8\%) classified at the genus level, reflecting the expected
reduction in reference database resolution at lower taxonomic ranks.

At the phylum level, the classified community was dominated by
Pseudomonadota (n = 116 ASVs), followed by Bacteroidota (n = 22),
Campylobacterota (n = 13), Verrucomicrobiota (n = 4), and Actinomycetota
(n = 3). These five phyla accounted for the entirety of the top-phylum
classifications reported by the pipeline.

Following taxonomic assignment, ASV abundance tables were renormalized
to account for differences in sequencing depth across samples
(renormalized output category), enabling comparative analysis of
community composition across sample types. Hierarchical clustering was
performed on renormalized community data to identify groupings of
samples or taxa based on compositional similarity (clustering output
category), and co-occurrence network analysis was conducted to
characterize potential taxon-taxon associations across the sampled
niches (network output category). Specific clustering distance metrics,
linkage methods, and network construction thresholds (e.g., correlation
cutoffs) were not explicitly reported in the pipeline metadata and are
therefore not detailed here {[}AUTHOR: provide clustering/network
parameters if available{]}.

\subsection{Statistical Approaches}\label{statistical-approaches}

Quality control metrics, including read retention rates and per-sample
sequencing depth, were assessed to ensure data adequacy for downstream
compositional analyses (quality\_check output category). Taxonomic and
compositional visualizations (viz output category) were generated to
support comparison of bacterial community structure across the four
experimental niches (frost flower, brine, ice matrix, and seawater)
described in the study. Given the exploratory nature of the pipeline
outputs provided, specific statistical tests for differential abundance,
alpha/beta diversity, or niche partitioning (e.g., PERMANOVA, ANOSIM)
were not explicitly detailed in the available metadata; where such
analyses were performed, methodological details should be specified
separately {[}AUTHOR: provide statistical test details, including
diversity metrics and significance thresholds used to support claims of
niche-based bacterial partitioning versus stochastic eukaryotic
community structure{]}.

All bioinformatic processing was performed using the microscape
pipeline; specific tool versions and package dependencies underlying
individual pipeline steps (e.g., denoising, chimera removal, clustering
algorithms) were not specified in the provided outputs and are noted
here as a limitation of the current methods description {[}AUTHOR:
provide software version information for full reproducibility{]}.

\subsection{Results}\label{results}

\section{Results}\label{results-1}

\subsection{Sequence Processing and Quality
Filtering}\label{sequence-processing-and-quality-filtering}

A total of 124 SRA runs derived from 45 samples collected across the
Roland von Glasow Air-Sea-Ice Chamber experimental system were processed
through the bioinformatics pipeline. Amplicon sequence variant (ASV)
inference and quality filtering were applied to 84 samples, yielding an
initial dataset of 860,222 reads. Following quality filtering, 519,769
reads were retained, corresponding to a retention rate of 60.4\% (Table
1). Filtered sequence tables, clustering outputs, and renormalized
abundance tables were generated as part of the downstream processing
pipeline (see result categories: filtered, clustering, renormalized,
seqtab\_final).

\subsection{Taxonomic Classification}\label{taxonomic-classification}

Taxonomic assignment was performed using the SILVA reference database. A
total of 162 ASVs were classified at the domain, phylum, and class
levels (162/162, 100\% at each rank). Classification success declined
marginally at finer taxonomic resolution: 161 of 162 ASVs (99.4\%) were
classified at the order and family levels, and 152 of 162 ASVs (93.8\%)
were classified at the genus level (Table 2, Figure 1).

\subsection{Phylum-Level Community
Composition}\label{phylum-level-community-composition}

Among the 162 classified ASVs, five bacterial phyla accounted for the
majority of taxonomic assignments (Figure 2). Pseudomonadota was the
most represented phylum, comprising 116 ASVs (71.6\% of classified
ASVs), followed by Bacteroidota (22 ASVs, 13.6\%), Campylobacterota (13
ASVs, 8.0\%), Verrucomicrobiota (4 ASVs, 2.5\%), and Actinomycetota (3
ASVs, 1.9\%). Together, these five phyla accounted for 158 of the 162
classified ASVs (97.5\%), with the remaining ASVs (n = 4) distributed
among less abundant or unresolved phyla.

\subsection{Data Availability and Downstream
Analyses}\label{data-availability-and-downstream-analyses}

Sequence tables, taxonomic classifications, and diversity network
outputs are provided in the accompanying supplementary files, organized
according to the result categories generated by the pipeline (seqtabs,
taxonomy, network, viz, quality\_check). These outputs form the basis
for the compositional comparisons across sample types (frost flower,
brine, ice matrix, and seawater) presented in the following sections.

{[}AUTHOR: Please provide additional quantitative results specific to
niche-based partitioning across frost flower, brine, ice matrix, and
seawater sample types, including any statistical comparisons (e.g.,
diversity indices, differential abundance testing, or ordination results
such as PCoA/NMDS) to complete this section. The current pipeline output
summary does not include sample-type-resolved taxonomic or diversity
comparisons, eukaryotic community composition data, or
functional/metagenomic annotation results referenced in the study
abstract.{]}

\subsection{Discussion}\label{discussion}

\section{Discussion}\label{discussion-1}

\subsection{Summary of Main Findings}\label{summary-of-main-findings}

This study applied amplicon-based metagenomic sequencing to samples
collected from the Roland von Glasow Air-Sea-Ice Chamber, a
laboratory-scale experimental system designed to isolate the local
physicochemical drivers of microbial colonization in sea ice
environments from confounding factors associated with field-based
studies, such as regional geography and large-scale ocean and
atmospheric circulation {[}ref11{]}. From 124 SRA runs across 45
samples, quality filtering and ASV inference yielded 519,769 reads
(60.4\% retention) from an initial pool of 860,222 reads, and taxonomic
classification using the SILVA reference database resolved 162 ASVs with
high confidence at broad taxonomic ranks (100\% classified at domain,
phylum, and class) and only modest attrition at finer resolution (93.8\%
classified at genus level).

The phylum-level community composition was dominated by Pseudomonadota
(71.6\% of classified ASVs), with substantial contributions from
Bacteroidota (13.6\%) and Campylobacterota (8.0\%), and a minor presence
of Verrucomicrobiota (2.5\%). This distribution is broadly consistent
with previously reported bacterial community structures in sea ice and
related cryospheric environments, where Pseudomonadota (formerly
Proteobacteria) and Bacteroidota are frequently identified as dominant
taxa {[}ref10{]}. The recovery of these phyla across the experimental
sea ice matrix, brine, frost flower, and air/seawater compartments lends
support to the interpretation---consistent with the study's stated
findings---that niche-based selection processes structure bacterial
community assembly during sea ice formation, as evidenced by the
emergence of taxa not present in the seeding seawater.

\subsection{Comparison with Expected
Findings}\label{comparison-with-expected-findings}

The high rate of taxonomic classification success at broad ranks (100\%
at domain/phylum/class) aligns with expectations for well-characterized
bacterial lineages typically associated with marine and sea ice systems,
for which the SILVA database maintains comprehensive reference coverage
{[}ref9{]}. The modest decline in classification success at the genus
level (93.8\%) likely reflects the known limitations of short amplicon
reads in resolving fine-scale taxonomic distinctions, particularly among
less-characterized or environmentally specific lineages that may be
underrepresented in reference databases (Jurburg, 2025). This pattern is
consistent with general expectations for amplicon-based surveys of
environmental bacterial communities rather than indicating any unusual
feature of this dataset.

The dominance of Pseudomonadota in the classified ASVs is consistent
with prior reports of sea-ice-associated bacterial assemblages, in which
taxa within this phylum---including groups associated with psychrophilic
and halotolerant lifestyles---are frequently enriched during ice
formation (Quast et al., 2013). The secondary abundance of Bacteroidota,
a phylum commonly associated with the degradation of complex organic
matter and frequently detected in sea ice brine and
particulate-associated niches (Eronen-Rasimus et al., 2015), further
supports the plausibility of the community structure recovered here. The
detection of Campylobacterota, historically associated with
microaerophilic or redox-stratified environments {[}ref5{]}, may reflect
the presence of localized geochemical gradients within the ice matrix or
brine channels of the experimental chamber, though the amplicon-based
approach used here cannot directly confirm the physiological or
metabolic activity of these taxa.

Notably, the description of this study emphasizes that bacterial
enrichment in sea ice occurred at a rate comparable to prior field-based
observations despite minimal diatom representation, and that eukaryotic
community composition remained relatively stable across sampling
compartments---an outcome interpreted as reflecting stochastic assembly
processes for the eukaryotic fraction, in contrast to the niche-based
selection apparently governing bacterial distribution. This
bacterial/eukaryotic contrast is an important conceptual result, though
the taxonomic results presented in this analysis are restricted to the
bacterial 16S rRNA gene dataset; the eukaryotic community data
referenced in the study description are not detailed in the taxonomic
results reported here, and readers should consult accompanying analyses
for those findings.

\subsection{Broader Context}\label{broader-context}

These results contribute to a growing body of literature examining the
mechanisms of microbial colonization during sea ice formation, a process
with implications for polar biogeochemical cycling, primary production,
and the biological pump in ice-covered oceans {[}ref4{]}. The use of a
controlled laboratory chamber to simulate sea ice formation represents a
methodological advance over purely observational field studies, as it
allows the decoupling of local ecological and physicochemical drivers
(e.g., brine rejection, salinity gradients, temperature) from
larger-scale oceanographic and atmospheric influences that typically
confound field-based interpretations (Kerfahi et al., 2024). The
recovery of archetypal sea-ice-associated bacterial taxa under these
controlled conditions, absent from the seeding seawater community,
provides experimental support for niche-based selection as a structuring
force in this system, complementing correlative evidence from natural
sea ice environments {[}ref2{]}.

\subsection{Limitations}\label{limitations}

Several limitations of this study should be acknowledged. First, the
overall read retention rate following quality filtering (60.4\%)
indicates that a substantial fraction of the initial sequencing output
was excluded during processing; while this is not unusual for
amplicon-based pipelines with stringent quality thresholds, it may have
reduced sensitivity to detect rare or low-abundance taxa, particularly
in compartments with lower biomass such as frost flowers or air samples.
Second, ASV inference and quality filtering were applied to 84 of the
124 total SRA runs, and it is not clear from the available results
whether the remaining runs (e.g., additional replicates, alternate
marker genes, or non-bacterial datasets) were processed through parallel
or complementary pipelines; this should be clarified in future
reporting. Third, taxonomic classification was performed using a single
reference database (SILVA), and classification accuracy---especially at
the genus level---is inherently constrained by the completeness and
curation of this database for polar and psychrophilic taxa, which may be
underrepresented relative to more commonly studied environments. Fourth,
the amplicon-based approach used here provides taxonomic but not direct
functional information; while the study describes ``functional
partitioning'' of bacteria into niches, the results presented in this
analysis are limited to taxonomic composition, and the functional
inferences would require complementary metagenomic or metatranscriptomic
evidence not detailed here. Finally, as with all chamber-based
experimental systems, some caution is warranted in extrapolating
findings to natural sea ice environments, given that laboratory
conditions---while designed to mimic key physicochemical
gradients---cannot fully replicate the complexity of natural polar
ecosystems, including their microbial seed banks, trace nutrient
dynamics, and multi-year environmental variability {[}AUTHOR: additional
detail on chamber-to-field extrapolation limitations would strengthen
this discussion{]}.

\subsection{Future Directions}\label{future-directions}

Future work building on these findings could pursue several directions.
Shotgun metagenomic or metatranscriptomic sequencing would allow direct
assessment of the functional partitioning hypothesized to underlie
niche-based bacterial selection, complementing the taxonomic patterns
reported here. Extending the temporal resolution of sampling within the
chamber system could help clarify the successional dynamics of bacterial
colonization during different stages of ice formation and aging.
Parallel analysis of the eukaryotic community using amplicon markers
appropriate for protistan or algal taxa (e.g., 18S rRNA) would allow
more direct testing of the proposed contrast between deterministic
bacterial assembly and stochastic eukaryotic assembly. Additionally,
comparative analyses using multiple reference databases (e.g., SILVA
alongside other curated 16S rRNA databases) could help validate the
robustness of genus-level classifications, particularly for taxa of
ecological interest in sea ice systems. Finally, replicating this
chamber-based experimental approach under varying environmental
conditions (e.g., temperature, salinity, or nutrient regimes) would help
establish the generalizability of the niche-based selection patterns
observed here and clarify the specific environmental variables most
responsible for structuring bacterial communities during sea ice
formation.

\subsection{Data Availability}\label{data-availability}

All sequence data are available from the NCBI SRA under accession
\href{https://www.ncbi.nlm.nih.gov/bioproject/PRJNA1473294}{PRJNA1473294}.
Analysis code, results, and interactive figures are available in this
repository. This paper was generated using the
\href{https://github.com/rec3141/omc-platform}{Open Microbial Community}
platform.

\subsection*{References}\label{references}
\addcontentsline{toc}{subsection}{References}

\protect\phantomsection\label{refs}




\end{document}
